\begin{definition}[Relative Minimum]
\label{def:relative-minimum}
Let $A \subseteq \mathbb{R}$, let $f : A \to \mathbb{R}$, and let $c \in A$. We say that $f$ has a \emph{relative minimum} at $c$ if there exists $\delta > 0$ such that
\[
f(c) \leq f(x)
\qquad \text{for all } x \in A \cap (c-\delta,c+\delta).
\]
If, in addition,
\[
f(c) < f(x)
\qquad \text{for all } x \in \bigl(A \cap (c-\delta,c+\delta)\bigr)\setminus\{c\},
\]
then $f$ has a \emph{strict relative minimum} at $c$.
\end{definition}

\begin{definition}[Relative Maximum]
\label{def:relative-maximum}
Let $A \subseteq \mathbb{R}$, let $f : A \to \mathbb{R}$, and let $c \in A$. We say that $f$ has a \emph{relative maximum} at $c$ if there exists $\delta > 0$ such that
\[
f(x) \leq f(c)
\qquad \text{for all } x \in A \cap (c-\delta,c+\delta).
\]
If, in addition,
\[
f(x) < f(c)
\qquad \text{for all } x \in \bigl(A \cap (c-\delta,c+\delta)\bigr)\setminus\{c\},
\]
then $f$ has a \emph{strict relative maximum} at $c$.
\end{definition}

\begin{definition}[Relative Extremum]
\label{def:relative-extremum}
Let $A \subseteq \mathbb{R}$, let $f : A \to \mathbb{R}$, and let $c \in A$. We say that $f$ has a \emph{relative extremum} at $c$ if $f$ has either a relative minimum at $c$ or a relative maximum at $c$.
\end{definition}

\begin{theorem}[Interior Extremum Theorem]
\label{thm:interior-extremum-theorem}
Let $I \subseteq \mathbb{R}$ be an interval, let $f : I \to \mathbb{R}$, and let $c \in \operatorname{int}(I)$. If $f$ has a relative minimum at $c$, then there exists $\delta > 0$ such that
\[
f(c) \leq f(x)
\qquad \text{for all } x \in I \cap (c-\delta,c+\delta).
\]
If $f$ has a relative maximum at $c$, then there exists $\delta > 0$ such that
\[
f(x) \leq f(c)
\qquad \text{for all } x \in I \cap (c-\delta,c+\delta).
\]
\end{theorem}

\begin{theorem}[Necessary Condition for an Interior Extremum]
\label{thm:necessary-condition-extremum}
Let $I \subseteq \mathbb{R}$ be an open interval, let $f : I \to \mathbb{R}$, and let $c \in I$. If $f$ has a relative extremum at $c$, then either $f'(c)$ does not exist, or
\[
f'(c)=0.
\]
\end{theorem}







\begin{theorem}[Rolle's Theorem]
\label{thm:rolles-theorem}
Let $I \subseteq \mathbb{R}$ be an interval, and let $a,b \in I$ with $a<b$. Suppose that $f : [a,b] \to \mathbb{R}$ is continuous on $[a,b]$, differentiable on $(a,b)$, and satisfies
\[
f(a)=f(b).
\]
Then there exists $c \in (a,b)$ such that
\[
f'(c)=0.
\]
\end{theorem}

\begin{theorem}[Mean Value Theorem]
\label{thm:mean-value-theorem}
Let $I \subseteq \mathbb{R}$ be an interval, and let $a,b \in I$ with $a<b$. Suppose that $f : [a,b] \to \mathbb{R}$ is continuous on $[a,b]$ and differentiable on $(a,b)$. Then there exists $c \in (a,b)$ such that
\[
f'(c)=\frac{f(b)-f(a)}{b-a}.
\]
\end{theorem}





\begin{remark*}[Common error]
The MVT guarantees \emph{existence} of $c$ but gives no formula for
it. The point $c$ depends on $f$, $a$, and $b$ and typically cannot
be identified explicitly. Using the MVT to compute a specific $c$ is
a misreading of the theorem — it is an existence result, not a
construction.
\end{remark*}

\begin{remark*}[Consequence]
The MVT is the chapter's load-bearing result. Its immediate corollaries
include: a function with $f' = 0$ on $(a,b)$ is constant
(\hyperref[cor:mvt-constant]{Constant Function Corollary}); a function
with $f' > 0$ is strictly increasing
(\hyperref[cor:mvt-monotonicity]{Monotonicity Corollary}); and a
function with $|f'| \le M$ is Lipschitz with constant $M$
(\hyperref[cor:mvt-lipschitz]{Lipschitz Bound Corollary}). Each of
these flows from the MVT by a one-line argument.
\end{remark*}





